%%%%%%%%%%%%%%%%%%%%%%%%%%%%%%%%%%%%%%%%%%%%%%%%%%%%%%%%%%%%%%
%%%%%%%%%%%%%%%%%% MA-0709 Geometría algebraica I %%%%%%%%%%%%%%%%%%%%%%%%%%%%%%
%%%%%%%%%%%%%%%%%%%%%%%%%%%%%%%%%%%%%%%%%%%%%%%%%%%%%%%%%%%%%%
\newpage
\subsection*{MA-0709 Geometría algebraica I}
\addcontentsline{toc}{subsection}{MA-0709 Geometría algebraica I}

\hypertarget{MA0709}{}
\begin{refsection}
\begin{table}[H]
\centering{
\begin{tabular}{|ll|}
\hline
%\multicolumn{2}{|c|}{\textbf{Nivel de virtualidad del curso:} Presencial}\\ \hline
\textbf{Año:} IV  & \textbf{Requisitos:} MA-0661 Álgebra Abstracta II \\
\textbf{Ciclo:} VIII  & \textbf{Correquisitos:} Ninguno \\
\textbf{Tipo de curso:} Teórico & \textbf{Horas presenciales por semana:} 5 \\
\textbf{Créditos:} 4 & \textbf{Horas de trabajo independiente por semana:} 7 \\
\textbf{Virtualidad:} P  &  \\ \hline
\end{tabular}
}
\end{table}

\subsubsection*{Descripción del curso}

Este curso introduce a la persona estudiante a las ideas y métodos fundamentales de la Geometría Algebraica mediante el estudio de conjuntos algebraicos, variedades y, especialmente, curvas algebraicas. El curso adopta un enfoque clásico, basado principalmente en el estudio de ecuaciones polinomiales y sus propiedades geométricas, siguiendo en buena medida el desarrollo del libro de William Fulton \cite{Ful89}.

Se comienza con conjuntos algebraicos afines, ideales y el Teorema de los Ceros de Hilbert, estableciendo la correspondencia fundamental entre objetos geométricos y algebraicos. Posteriormente se estudian variedades afines, anillos de coordenadas, funciones racionales y propiedades locales. Las curvas algebraicas planas proporcionan el contexto principal para introducir multiplicidades, tangentes, puntos singulares y puntos suaves, así como números de intersección.

Una segunda parte del curso desarrolla la geometría proyectiva y las curvas proyectivas planas, incluyendo el teorema de Bézout y resultados clásicos sobre intersecciones. Se estudian también morfismos, aplicaciones racionales y equivalencia birracional. Estas herramientas permiten abordar la resolución de singularidades de curvas mediante transformaciones cuadráticas y explosiones.

El estudio de curvas proyectivas no singulares conduce a la teoría de divisores, diferenciales, género y al Teorema de Riemann--Roch, uno de los resultados centrales del curso. Finalmente, se introduce la geometría algebraica sobre cuerpos finitos, la función zeta de una variedad algebraica y las conjeturas de Weil, acompañadas de ejemplos explícitos de conteo de puntos y cálculo de funciones zeta.

El curso busca mantener un equilibrio entre el desarrollo teórico y el estudio detallado de ejemplos concretos de curvas. Cuando resulte apropiado, se podrán utilizar sistemas de álgebra computacional para estudiar ecuaciones, singularidades, intersecciones y puntos sobre cuerpos finitos.

\subsubsection*{Objetivos}

\textbf{General:} Introducir a la persona estudiante a los conceptos, métodos y resultados fundamentales de la Geometría Algebraica clásica, con énfasis en variedades y curvas algebraicas, sus singularidades, propiedades proyectivas y birracionales, y sus invariantes geométricos y aritméticos. \

\textbf{Específicos:}

\begin{enumerate}
\item Comprender la relación entre conjuntos algebraicos e ideales de anillos de polinomios.
\item Aplicar el Teorema de los Ceros de Hilbert al estudio de variedades algebraicas afines.
\item Estudiar propiedades locales de curvas algebraicas, incluyendo multiplicidades, tangentes, singularidades y puntos suaves.
\item Comprender los fundamentos de la geometría proyectiva y su aplicación al estudio de curvas algebraicas.
\item Analizar intersecciones de curvas y aplicar el teorema de Bézout en ejemplos concretos.
\item Comprender los conceptos de morfismo, aplicación racional y equivalencia birracional.
\item Estudiar la resolución de singularidades de curvas algebraicas.
\item Introducir la teoría de divisores y comprender el género y el Teorema de Riemann--Roch para curvas.
\item Estudiar curvas y variedades sobre cuerpos finitos y calcular ejemplos de sus funciones zeta.
\item Conocer el enunciado y significado de las conjeturas de Weil y su relación con la geometría y la aritmética.
\item Comunicar argumentos, cálculos y resultados de manera clara y rigurosa en forma oral y escrita.
\end{enumerate}

\subsubsection*{Contenidos}

\begin{enumerate}

\item \textbf{Conjuntos algebraicos afines y el Teorema de los Ceros de Hilbert (2 semanas):}
\begin{enumerate}
\item Espacio afín y conjuntos algebraicos.
\item Ideales asociados a conjuntos algebraicos y conjuntos de ceros asociados a ideales.
\item Componentes irreducibles y descomposición de conjuntos algebraicos.
\item Teorema de la base de Hilbert.
\item Teorema de los Ceros de Hilbert y correspondencia entre geometría y álgebra.
\end{enumerate}

\item \textbf{Variedades afines y propiedades locales (2 semanas):}
\begin{enumerate}
\item Variedades afines y anillos de coordenadas.
\item Aplicaciones polinomiales y cambios de coordenadas.
\item Funciones racionales y cuerpos de funciones.
\item Anillos locales y su interpretación geométrica.
\item Dimensión y primeros ejemplos.
\end{enumerate}

\item \textbf{Propiedades locales de curvas planas (2 semanas):}
\begin{enumerate}
\item Curvas algebraicas planas y sus componentes.
\item Multiplicidad de un punto y cono tangente.
\item Puntos suaves y puntos singulares.
\item Tangentes y singularidades clásicas: nodos, cúspides y puntos múltiples.
\item Multiplicidad de intersección de curvas.
\end{enumerate}

\item \textbf{Geometría proyectiva (2 semanas):}
\begin{enumerate}
\item Espacio proyectivo y coordenadas homogéneas.
\item Polinomios homogéneos y conjuntos algebraicos proyectivos.
\item Clausura proyectiva de variedades afines.
\item Variedades proyectivas e hiperplanos.
\item Relación entre geometría afín y proyectiva.
\end{enumerate}

\item \textbf{Curvas proyectivas planas e intersecciones (2 semanas):}
\begin{enumerate}
\item Curvas proyectivas planas.
\item Sistemas lineales de curvas.
\item Teorema de Bézout.
\item Puntos múltiples e intersecciones.
\item Teorema fundamental de Max Noether y aplicaciones seleccionadas.
\end{enumerate}

\item \textbf{Morfismos y geometría birracional (1 semana):}
\begin{enumerate}
\item Variedades y morfismos.
\item Funciones y aplicaciones racionales.
\item Cuerpos de funciones.
\item Equivalencia birracional de curvas.
\item Modelos afines y proyectivos de una misma curva.
\end{enumerate}

\item \textbf{Resolución de singularidades de curvas (1 semana):}
\begin{enumerate}
\item Transformaciones birracionales de curvas.
\item Explosión de un punto y transformaciones cuadráticas.
\item Transformadas propias de curvas.
\item Resolución de singularidades de curvas planas.
\item Ejemplos explícitos de resolución de nodos, cúspides y otras singularidades sencillas.
\end{enumerate}

\item \textbf{Divisores, género y Teorema de Riemann--Roch (2 semanas):}
\begin{enumerate}
\item Curvas proyectivas no singulares y sus cuerpos de funciones.
\item Divisores y divisores principales.
\item Espacios de funciones asociados a divisores.
\item Diferenciales y divisores canónicos.
\item Género de una curva.
\item Teorema de Riemann y Teorema de Riemann--Roch.
\item Consecuencias y ejemplos en curvas de género bajo.
\end{enumerate}

\item \textbf{Curvas sobre cuerpos finitos, funciones zeta y conjeturas de Weil (1 semana):}
\begin{enumerate}
\item Variedades algebraicas sobre cuerpos finitos y conteo de puntos sobre extensiones finitas.
\item Definición de la función zeta de una variedad algebraica sobre un cuerpo finito.
\item Cálculos explícitos para espacios afines, espacios proyectivos y curvas sencillas.
\item Función zeta de curvas proyectivas no singulares.
\item Ejemplos de curvas elípticas sobre cuerpos finitos.
\item Enunciado y significado de las conjeturas de Weil.
\item Racionalidad, ecuación funcional y análogo de la Hipótesis de Riemann sobre cuerpos finitos, a nivel de enunciados.
\end{enumerate}

\end{enumerate}

\subsubsection*{Metodología}

\noindent \textbf{Lineamientos metodológicos:} La persona docente a cargo de este curso debe respetar las siguientes pautas:
\begin{enumerate}
\item El curso debe desarrollarse principalmente desde el punto de vista clásico de variedades y curvas algebraicas, sin requerir el lenguaje general de esquemas.
\item Los conceptos algebraicos deben relacionarse sistemáticamente con su interpretación geométrica.
\item El estudio de curvas algebraicas debe ocupar un lugar central y servir como fuente constante de ejemplos concretos.
\item Las singularidades, las intersecciones y las transformaciones birracionales deben ilustrarse mediante ejemplos explícitos.
\item El Teorema de Riemann--Roch debe presentarse como uno de los resultados centrales que integran los distintos temas desarrollados a lo largo del curso.
\item La sección final sobre funciones zeta y conjeturas de Weil debe enfatizar la conexión entre geometría algebraica y teoría de números.
\end{enumerate}

\textbf{Sugerencias metodológicas:} Adicionalmente, se tienen las siguientes recomendaciones para la persona docente a cargo de este curso:
\begin{enumerate}
\item Utilizar desde el inicio ejemplos concretos de cónicas, cúbicas y curvas de grados bajos para desarrollar intuición geométrica.
\item Complementar las definiciones de singularidad, multiplicidad e intersección con representaciones gráficas siempre que sea posible.
\item Utilizar herramientas computacionales como \texttt{SageMath}, \texttt{Singular} u otros sistemas de álgebra computacional para estudiar ideales, intersecciones, singularidades y ejemplos de curvas.
\item Realizar cálculos explícitos de explosiones y resolución de singularidades en ejemplos sencillos.
\item Complementar el estudio del Teorema de Riemann--Roch con ejemplos en la recta proyectiva, cónicas, curvas cúbicas y otras curvas de género pequeño.
\item Utilizar experimentación computacional para contar puntos de curvas sobre cuerpos finitos y calcular los primeros términos de sus funciones zeta.
\item Promover la resolución de problemas como parte central del aprendizaje, utilizando en particular referencias orientadas al aprendizaje mediante ejercicios.
\item Relacionar cuando sea apropiado los temas estudiados con teoría de números, superficies de Riemann, curvas elípticas y geometría aritmética.
\item En la evaluación, se sugiere combinar problemas teóricos, cálculos explícitos y pequeños proyectos o exposiciones.
\item Promover la comunicación clara y rigurosa de argumentos e ideas matemáticas de manera oral y escrita.
\end{enumerate}

\subsubsection*{Guía para las referencias}

El libro de William Fulton \cite{Ful89} constituye la referencia principal del curso. Su desarrollo parte de conjuntos y variedades algebraicas afines, continúa con las propiedades locales de curvas planas y la geometría proyectiva, y posteriormente estudia curvas proyectivas planas, morfismos, aplicaciones racionales, resolución de singularidades y el Teorema de Riemann--Roch. Esta secuencia determina en buena medida la organización del presente curso.

El libro \emph{Algebraic Geometry: A Problem Solving Approach} de Garrity y colaboradores \cite{Gar13} es un complemento especialmente útil por su enfoque basado en problemas. Puede utilizarse como fuente de ejercicios y actividades para desarrollar intuición geométrica. Los textos de Smith, Kahanpää, Kekäläinen y Traves \cite{SKKT00}, Hulek \cite{Hul03} y Holme \cite{Hol12} proporcionan introducciones alternativas a variedades afines y proyectivas, suavidad, geometría birracional y teoría de curvas.

Los libros de Ernst Kunz \cite{Kun13} y \cite{Kun05} son particularmente útiles para estudiar la relación entre álgebra conmutativa y geometría algebraica, así como propiedades locales y singularidades de curvas. Para un tratamiento más geométrico y concreto de curvas planas se recomiendan también los libros de Fischer \cite{Fis01} y Kendig \cite{Ken11}. El texto de Kirwan \cite{Kir92} ofrece una perspectiva complementaria sobre curvas algebraicas complejas, superficies de Riemann, diferenciales, género y singularidades. Por su énfasis en curvas algebraicas y sus aspectos analíticos y topológicos, es particularmente apropiado como lectura complementaria en la parte final del curso.

El libro de Ueno \cite{Uen97} es especialmente útil para el enfoque de este curso: además de estudiar variedades proyectivas, curvas y resolución de singularidades, contiene un tratamiento de las funciones zeta de congruencias de curvas algebraicas sobre cuerpos finitos. El libro de Bump \cite{Bum98} constituye otra referencia importante para la última parte del curso. Desarrolla curvas proyectivas no singulares, diferenciales, el Teorema de Riemann--Roch y culmina con un capítulo dedicado a la función zeta de una curva sobre un cuerpo finito, incluyendo la Hipótesis de Riemann para curvas elípticas. Por tanto, \cite{Uen97} y \cite{Bum98} son las referencias recomendadas específicamente para el tema de funciones zeta y su relación con las conjeturas de Weil.

\nocite{Ful89}
\nocite{Gar13}
\nocite{Uen97}
\nocite{SKKT00}
\nocite{Kun13}
\nocite{Hol12}
\nocite{Bum98}
\nocite{Hul03}
\nocite{Kun05}
\nocite{Fis01}
\nocite{Ken11}
\nocite{Kir92}

\printcursosbibliography
\end{refsection}
